\begin{textsample}{POS Dim 3 – es – Score 19.00 – es/es\_inf\_8\_camp\_exp.txt}  \label{ex:f3_pos_119}
Campo de Experiências : O Eu, o Outro e o Nós A educação busca garantir o direito individual ao pleno desenvolvimento do potencial de cada pessoa, como um passo fundamental para que todos possam alcançar uma vida valiosa, de sua própria escolha.
Contudo, para que a \textbf{criança} alcance o pleno desenvolvimento, além de adquirir as competências e o conhecimento sobre si mesma, sobre o outro e sobre o mundo, precisa que lhe sejam dadas condições e oportunidades efetivas para aplicação das competências e dos conhecimentos \textbf{assimilados} no seu cotidiano.
No convívio com o outro, a \textbf{criança} se constitui enquanto sujeito com um modo próprio de pensar, agir e sentir.
É nessa interação que ela constrói sua identidade.
Esse processo acontece ao longo da vida, mas a \textbf{criança} o vive de forma intensa na primeira \textbf{infância}.
Por isso, é de grande importância oportunizar a ela, nos espaços da Educação \textbf{Infantil}, novas formas de conviver que ampliem sua confiança e participação nas relações que estabelece com o outro.
É preciso oportunizá -la a se conhecer como alguém que tem características próprias, concepções ( apesar de muito \textbf{pequena} ainda ), \textbf{desejos}, motivações e que se interrelaciona com o outro que também tem \textbf{desejos} e interesses próprios, se conscientizando da existência de um “ nós ” enquanto seres dependentes uns dos outros que nos constituímos nessa relação ampliada e diversa.
O foco deste Campo de Experiências é proporcionar à \textbf{criança} vivenciar diferentes situações de atenção pessoal e outras práticas sociais, formas mais democráticas, respeitosas, de cooperação e solidariedade no relacionamento com seus \textbf{pares} e \textbf{adultos}.
É desafiador para a \textbf{criança} perceber essas diferenças e compreender que as pessoas exercem diferentes papéis em relação ao eu e compreender que as culturas, as formas de linguagem e a constituição familiar se diferenciam nos modos de viver.
Quando passa a frequentar a Educação \textbf{Infantil}, a \textbf{criança} é \textbf{acolhida}, interage com outros \textbf{parceiros}, cria novos \textbf{laços} afetivos de convivência, expressa suas emoções, pensamentos, sentimentos e \textbf{percepções} e confronta suas formas de viver com a desses \textbf{parceiros}, construindo uma identidade livre de preconceitos, de raça, cor, religião, condição social, entre outros, ampliando suas possibilidades de \textbf{cuidar} de si e do outro.
Na BNCC este Campo de Experiências estabelece que : > É na interação com os \textbf{pares} e com \textbf{adultos} que as \textbf{crianças} vão constituindo um modo próprio de agir, sentir e pensar e vão descobrindo que existem outros modos de vida, pessoas diferentes, com outros pontos de vista.
Conforme vivem suas primeiras experiências sociais ( na família, na \textbf{instituição} escolar, na coletividade ), constroem \textbf{percepções} e questionamentos sobre si e sobre os outros, diferenciando -se e, simultaneamente, identificando -se como seres individuais e sociais.
Ao mesmo tempo em que participam de relações sociais e de \textbf{cuidados} pessoais, as \textbf{crianças} constroem sua autonomia e senso de autocuidado, de reciprocidade e de interdependência com o meio.
Por sua vez, na Educação \textbf{Infantil}, é preciso criar oportunidades para que as \textbf{crianças} entrem em contato com outros grupos sociais e culturais, outros modos de vida, diferentes atitudes, técnicas e rituais de \textbf{cuidados} pessoais e do grupo, costumes, celebrações e narrativas.
Nessas experiências, elas podem ampliar o modo de perceber a si mesmas e ao outro, valorizar sua identidade, respeitar os outros e reconhecer as diferenças que nos constituem como \textbf{seres} humanos. ( BRASIL, 2017, p.36 ) Assim, a ênfase do Campo O Eu, o Outro e o Nós está ligada a constituição de atitudes nas relações vivenciadas pela \textbf{criança} ao longo da Educação \textbf{Infantil}, colocando as interações e \textbf{brincadeiras} como eixos do processo educativo e tratando dos Direitos de Aprendizagem que entrelaçam as experiências concretas da vida cotidiana das \textbf{crianças} com os conhecimentos sistematizados possibilitando à esta : Figura 22 - Direitos de Aprendizagem aplicados ao Campo “ O Eu, o Outro e o Nós ” ( \textbf{OLIVEIRA}, 2017, p.22 ).
Conviver - com \textbf{crianças} e \textbf{adultos} em \textbf{pequenos} grupos, reconhecendo e respeitando as diferentes identidades e pertencimento étnico-racial, de gênero e religião de seus \textbf{parceiros}. \textbf{Brincar} - com diferentes \textbf{parceiros}, desenvolvendo sua imaginação e solidariedade.
Participar - ativamente das situações do cotidiano, tanto daquelas ligadas ao \textbf{cuidado} de si e do ambiente, como das relativas às atividades propostas pelo/a professor/a.
Explorar - diferentes formas de interagir com \textbf{parceiros} diversos em situações variadas, ampliando sua noção de mundo e sua sensibilidade em relação aos outros.
Expressar - às outras \textbf{crianças} e/ou \textbf{adultos} suas necessidades, emoções, sentimentos, dúvidas, hipóteses, descobertas, opiniões, oposições.
Conhecer -se - e construir uma identidade pessoal e cultural, valorizando suas características e as das outras \textbf{crianças} e \textbf{adultos}, aprendendo a identificar e combater atitudes preconceituosas e discriminatórias.

% matched lemmas: acolher, adulto, assimilar, brincadeira, brincar, criança, cuidado, cuidar, desejo, infantil, infância, instituição, laço, oliveira, par, parceiro, pequeno, percepção, seres
\end{textsample}
